%!TEX program = xelatex
%!TEX encoding = UTF-8
\def\ano{2026}
\def\semestre{1}
\def\disciplina{Geometria Analítica}
\def\turma{1}
\def\autor{José Antônio O. Freitas}
\def\instituto{MAT-UnB}
\documentclass{beamer}
\usetheme{Madrid}
\usecolortheme{beaver}
% \mode<presentation>
\usepackage[singlelinecheck=false]{caption}
\usepackage{lmodern}
\usepackage{textpos}
\usepackage{amssymb}
\usepackage{amsmath,amsfonts,amsthm,amstext}
\usepackage{cases}
\usepackage[justification=centering]{caption}
\usepackage[brazilian]{babel}
% \usepackage[latin1]{inputenc}
\usepackage{graphicx}
\graphicspath{{/home/jfreitas/GitHub/geometria_analitica/apresentacoes/logo}{/home/jfreitas/GitHub/geometria_analitica/notas_de_aula/figuras-latex}{C:/Users/jfreitas/GitHub/geometria_analitica/apresentacoes/logo}}
\usepackage{enumitem}
\usepackage{multicol}
\usepackage{answers}
\usepackage{qrcode}
\usepackage{tikz,ifthen}
\usetikzlibrary{lindenmayersystems}
\usetikzlibrary[shadings]

\newcounter{saveenumi}
\newcommand{\seti}{\setcounter{saveenumi}{\value{enumi}}}
\newcommand{\conti}{\setcounter{enumi}{\value{saveenumi}}}

\resetcounteronoverlays{saveenumi}


\renewcommand{\qedsymbol}{$\blacksquare$}

\setbeamertemplate{itemize items}[ball]

\input{../configuracoes_latex/ambientes}
\newcommand{\nsub}{\varsubsetneq}
\newcommand{\vaz}{\emptyset}
\newcommand{\dom}{{\rm dom\,}}
\newcommand{\im}{{\rm Im\,}}
\newcommand{\sub}{\subseteq}
\newcommand{\n}{\mathbb{N}}
\newcommand{\z}{\mathbb{Z}}
\newcommand{\rac}{\mathbb{Q}}
\newcommand{\real}{\mathbb{R}}
\newcommand{\complex}{\mathbb{C}}
\newcommand{\cp}[1]{\mathbb{#1}}
\newcommand{\ch}{\mbox{\textrm{car\,}}\nobreak}
\newcommand{\vesp}[1]{\vspace{ #1  cm}}
\newcommand{\compcent}[1]{\vcenter{\hbox{$#1\circ$}}}
\newcommand{\comp}{\mathbin{\mathchoice
{\compcent\scriptstyle}{\compcent\scriptstyle}
{\compcent\scriptscriptstyle}{\compcent\scriptscriptstyle}}}
\newcommand\T{\rule{0pt}{2.6ex}}
\newcommand{\elipse}[2]{\dfrac{x^2}{#1} + \dfrac{y^2}{#2} = 1}
\newcommand{\inner}[2]{\langle \vec{#1}, \vec{#2} \rangle}
\newcommand{\proj}{{\rm proj}}
\newcommand{\eq}{=}
\newcommand{\norm}[1]{\left\Vert{#1}\right\Vert}
\newcommand{\integer}{\mathbb{Z}}
\renewcommand{\sin}{{\rm sen\,}}
\renewcommand{\tan}{{\rm tg\,}}
\renewcommand{\csc}{{\rm cossec\,}}
\renewcommand{\cot}{{\rm cotg\,}}
\renewcommand{\sinh}{{\rm senh\,}}




\title{Vetores no Espaço}
\author[\autor]{\autor}
\institute[\instituto]{\instituto}
\date{}

\begin{document}
\begin{frame}
    \maketitle
\end{frame}

\logo{
\includegraphics[scale=.1]{logo-MAT.png}\vspace*{8.5cm}}

\begin{frame}
    Fixemos um ponto $O$ no espaço, \pause que será denominado como \textbf{origem}. \pause Tomemos três retas \pause duas a duas perpendiculares entre si \pause e concorrentes em $O$, \pause que serão denominados \textbf{eixos coordenados} \pause e serão denotados por $Ox$, \pause $Oy$ \pause e $Oz$.
\end{frame}

\begin{frame}
    \begin{figure}[!h]
        \centering
        \caption{Sistema de coordenadas tridimensionais}
        \includegraphics{eixos-coordenados-espaco.pdf}
    \end{figure}
\end{frame}

\begin{frame}
    Projetando um ponto $P$ \pause do espaço ortogonalmente sobre cada um dos eixos coordenados \pause podemos representar os pontos do espaço por ternas ordenadas $(a,b,c)$ \pause de números reais.
\end{frame}

\begin{frame}
    \begin{figure}[!h]
        \centering
        \caption{Coordenadas de um ponto no espaço tridimensional}
        \includegraphics{coordenadas-ponto-espaco.pdf}
    \end{figure}
\end{frame}

\begin{frame}
    Assim os pontos de $\real^3$ \pause são descritos por\pause
    \[
        \real^3 = \{(a,b,c)\pause \mid a,b,c \in \real\}.
    \]
\end{frame}

\begin{frame}
    Sejam $P(x_1,y_1,z_1)$ \pause e $Q(x_2,y_2,z_2)$ \pause dois pontos de $\real^3$. \pause Traçando por $P$ um segmento $PS$ \pause paralelo ao segmento $P_0Q_0$, \pause onde $P_0(x_1,y_1)$ e $Q_0(x_2,y_2)$, \pause obtemos que\pause
    \[
        d(P,Q)^2 = d(P_0,Q_0)^2 + |SQ|^2.
    \]
\end{frame}

\begin{frame}
    \begin{figure}
        \centering
        \caption{Distância entre pontos em $\real^3$}
        \includegraphics{distancia-pontos-espaco.pdf}
    \end{figure}
\end{frame}

\begin{frame}
    Agora, \pause $|SQ| = |z_2 - z_1|$ \pause e daí\pause
    \[
        D(P,Q) = \sqrt{(x_1 - x_2)^2 + (y_1 - y_2)^2 + (z_1 - z_2)^2}.
    \]
\end{frame}

\begin{frame}
    Os vetores e as operações \pause com eles podem ser definidas de forma semelhante ao feito em $\real^2$. \pause Para isso, um ponto $P$ \pause do espaço é representado em $\real^3$ \pause por uma terna de coordenadas $P(x_1, y_1,z_1)$, \pause onde $x_1$, $y_1$, $z_1 \in \real$.\pause

    Assim escrevemos
    \[
        \vec{u} = (x_1,y_1,z_1).
    \]\pause
    para representar um vetor com extremidade inicial no ponto $(0, 0, 0)$ \pause e extremidade final no ponto $(x_1, y_1, z_1)$.\pause

    O vetor nulo $\vec{0}$ \pause é escrito como\pause
    \[
        \vec{0} = (0,0,0).
    \]
\end{frame}

\begin{frame}
    Assim temos:\pause
    \begin{definicao}
        Sejam $\alpha \in \real$, \pause $\vec{u} = (x_1, y_1,z_1)$ \pause e $\vec{v} = (x_2, y_2,z_2)$. \pause Então:\pause
        \begin{enumerate}[label=({\arabic*})]
            \item $\vec{u} + \vec{v} = (x_1 + x_2, y_1 + y_2,z_1 + z_2)$;\pause
            \item $\alpha \vec{u} = (\alpha x_1, \alpha y_1, \alpha z_1)$.
        \end{enumerate}
    \end{definicao}
\end{frame}

\begin{frame}
    Em alguns casos pode ser mais conveniente representar um vetor $\vec{u} = (x_1, y_1,z_1)$ \pause de forma matricial. \pause Para isso escrevemos as componentes de $\vec{u}$ \pause como uma matriz de 1 \pause coluna e três linhas\pause
    \[
        \vec{u} = \begin{bmatrix}
            x_1 \\y_1\\z_1
        \end{bmatrix}.
    \]
\end{frame}

\begin{frame}
    Vimos que a norma de um vetor em $\real^2$ \pause é definida como o comprimento de um segmento orientado que o represente. \pause Com a representação de vetores por coordenadas podemos reescrever o conceito de norma da seguinte maneira:\pause
    \begin{definicao}
        Seja $\vec{u} = (x_1, y_1,z_1)$ \pause um vetor. A \textbf{norma} \pause de $\vec{u}$, \pause denotada por $\norm{\vec{u}}$ \pause é dada por\pause
        \[
            \norm{\vec{u}} = \sqrt{x_1^2 + y_1^2 + z_1^2}.
        \]
    \end{definicao}
\end{frame}

\begin{frame}
    Um vetor de norma 1 \pause é chamado de \textbf{vetor unitário}. \pause Dado um vetor não nulo $\vec{u}$ \pause o vetor\pause
    \[
        \vec{v} = \dfrac{\vec{u}}{\norm{\vec{u}}}
    \]\pause
    é um vetor unitário \pause na direção de $\vec{u}$.
\end{frame}

\begin{frame}
    \begin{proposicao}
        Sejam $\vec{u} = (x_1, y_1,z_1)$ \pause um vetor em $\real^3$ \pause e $\alpha \in \real$. \pause Então:\pause
        \begin{enumerate}[label=({\arabic*})]
            \item Se $\vec{u} \ne \vec{0}$, \pause então $\norm{\vec{u}} \ne 0$.\pause
            \item $\norm{\vec{u}} \ne 0$ \pause se, e somente se, \pause $\vec{u} \ne \vec{0}$.\pause
            \item $\norm{\alpha\vec{u}} = \pause |\alpha|\norm{\vec{u}}$.
        \end{enumerate}
    \end{proposicao}
\end{frame}

\begin{frame}
    \begin{definicao}
        Sejam $\vec{u}$ \pause e $\vec{v}$ \pause vetores não nulos. \pause O \textbf{ângulo} \pause entre $\vec{u}$ e $\vec{v}$ \pause é definido pelo ângulo $\theta$ \pause determinado por $\vec{u}$ e $\vec{v}$ \pause e que satisfaz $0 \le \theta \le \pi$, \pause quando os vetores $\vec{u}$ e $\vec{v}$ \pause são representados com a mesma origem.
    \end{definicao}
\end{frame}

\begin{frame}
    \begin{definicao}
        Quando o ângulo $\theta$ \pause entre dois vetores $\vec{u}$ e $\vec{v}$ \pause é reto, isto é, $\theta = \pi/2$, \pause ou um deles é o vetor nulo, \pause dizemos que os vetores $\vec{u}$ e $\vec{v}$ são \textbf{ortogonais} \pause ou \textbf{perpendiculares} entre si. \pause Denotamos tal fato, escrevendo $\vec{u} \perp \vec{v}$.
    \end{definicao}
\end{frame}

\begin{frame}
    \begin{definicao}\label{produtointerno-espaco}
        O \textbf{produto escalar} \pause ou \textbf{produto interno} \pause dos vetores $\vec{u}$ \pause e $\vec{v}$, \pause indicado por $\inner{u}{v}$, \pause é o número real tal que\pause
        \begin{enumerate}[label=({\arabic*})]
            \item Se $\vec{u}$ \pause ou $\vec{v}$ é nulo, \pause então $\inner{u}{v} = 0$.\pause
            \item Se $\vec{u}$ e $\vec{v}$ \pause não são nulos \pause e $\theta$ é o ângulo entre $\vec{u}$ e $\vec{v}$, \pause então
                  \begin{align}\label{produto-interno-espaco}
                      \inner{u}{v} = \norm{\vec{u}}\norm{\vec{v}}\cos\theta.
                  \end{align}
        \end{enumerate}
    \end{definicao}
\end{frame}

\begin{frame}
    \begin{proposicao}
        Sejam $\vec{u}$ e $\vec{v}$ vetores \pause e $\theta$ o \pause ângulo entre $\vec{u}$ e $\vec{v}$.\pause
        \begin{enumerate}[label=({\arabic*})]
            \item Se $\vec{u}$ e $\vec{v}$ \pause não são nulos, \pause então\pause
                  \[
                      \cos\theta = \dfrac{\inner{u}{v}}{\norm{\vec{u}}\norm{\vec{v}}}.
                  \]\pause
            \item Qualquer que seja o vetor $\vec{u}$,\pause
                  \[
                      \norm{\vec{u}} = \sqrt{\inner{u}{u}}.
                  \]\pause
            \item Quaisquer que sejam os vetores $\vec{u}$ \pause e $\vec{v}$, \pause $\vec{u}\perp\vec{v}$ \pause se, e somente se, \pause $\inner{u}{v} = 0$.
        \end{enumerate}
    \end{proposicao}
\end{frame}

\begin{frame}
    De modo análogo ao caso para vetores no plano, \pause temos:\pause
    \begin{teorema}
        O produto interno de dois vetores \pause $\vec{u} = (x_1, y_1,z_1)$ \pause e $\vec{v} = (x_2, y_2,z_2)$ \pause em $\real^3$ \pause é dado por\pause
        \[
            \inner{u}{v} = \pause x_1x_2 \pause + y_1y_2 \pause + z_1z_2.
        \]
    \end{teorema}
\end{frame}

\begin{frame}
    \begin{proposicao}\label{propriedades-produto-vetorial-espaco}
        Quaisquer que sejam os vetores \pause $\vec{u} = (x_1, y_1,z_1)$, \pause $\vec{v} = (x_2, y_2,z_2)$ \pause e $\vec{w} = (x_3, y_3,z_3)$ \pause e qualquer que seja o número real $\lambda$ \pause temos\pause
        \begin{enumerate}[label=({\roman*})]
            \item\label{linearidade-soma-produto-vetorial-espaco} $\langle\vec{u}, \pause \vec{v} + \vec{w}\rangle \pause = \inner{u}{v} \pause + \inner{u}{w}$\pause
            \item\label{linearidade-escalar-produto-vetorial-espaco} $\langle\vec{u}, \pause \lambda\vec{v}\rangle \pause = \langle\lambda\vec{u}, \pause \vec{v}\rangle \pause = \lambda \pause \inner{u}{v}$\pause
            \item $\inner{u}{v} \pause = \inner{v}{u}$\pause
            \item Se $\vec{u} \ne \vec{0}$, \pause então $\inner{u}{u} > 0$.
        \end{enumerate}
    \end{proposicao}
\end{frame}

\begin{frame}
    \begin{observacao}
        \begin{enumerate}[label=({\alph*})]
            \item As propriedades \ref{linearidade-soma-produto-vetorial-espaco} \pause e \ref{linearidade-escalar-produto-vetorial-espaco} \pause da Proposição \ref{propriedades-produto-vetorial-espaco} \pause podem ser estendidas para qualquer número de vetores:\pause
                  \[
                      \langle\vec{u}, \pause \lambda_1\vec{v_1} \pause + \lambda_2\vec{v_2} \pause + \cdots \pause + \lambda_n\vec{v_n}\rangle \pause = \lambda_1\inner{u}{v_1} \pause + \lambda_2\inner{u}{v_2} \pause + \cdots \pause + \lambda_n\inner{u}{v_n}.\pause
                  \]
            \item Na igualdade $\inner{u}{v} = \inner{u}{w}$ \pause não podemos concluir que $\vec{v} = \vec{w}$. \pause Por exemplo, \pause para $\vec{u} = (1, 0, 0)$, \pause $\vec{v} = (2, 1, 1)$ e $\vec{w} = (2, -5, 3)$ \pause temos $\inner{u}{v} = \inner{u}{w}$ \pause e no entanto $\vec{v} \ne \vec{w}$. \pause Mas podemos concluir que \pause $\vec{u}\perp(\vec{v} - \vec{w})$.\pause

            \item De $\inner{u}{v} = 0$ \pause não podemos concluir que $\vec{u} = \vec{0}$ \pause ou $\vec{v} = \vec{0}$. \pause Por exemplo, \pause para $\vec{u} = (2, 1, 1)$ \pause e $\vec{v} = (-2, 4, 0)$ \pause temos $\inner{u}{v} = 0$ \pause e no entanto $\vec{u} \ne \vec{0}$ \pause e $\vec{v} \ne \vec{0}$.
        \end{enumerate}
    \end{observacao}
\end{frame}

\begin{frame}
    \begin{proposicao}
        Sejam $\vec{u}$ e $\vec{v}$ vetores.\pause
        \begin{enumerate}[label=({\arabic*})]
            \item $\mid\inner{u}{v}\mid \pause \le \norm{\vec{u}}\pause \norm{\vec{v}}$ \pause (Desigualdade de Scharwz)\pause
            \item $\norm{\vec{u} + \vec{v}} \pause \le \norm{\vec{u}} \pause + \norm{\vec{v}}$ \pause (Desigualdade Triangular)\pause
        \end{enumerate}
    \end{proposicao}
\end{frame}

\begin{frame}
    Dados vetores $\vec{u} = (a_1,b_1,c_1)$ \pause e $\vec{v} = (a_2,b_2,c_2)$ \pause é possível encontrar um vetor $\vec{w} = (x,y,z)$ \pause que seja simultaneamente \pause ortogonal a $\vec{u}$ \pause e a $\vec{v}$?
\end{frame}

\begin{frame}
    O vetor $\vec{w}$ dado por
    \begin{equation}\label{produtovetorial}
        \vec{w} = (b_1c_2 - b_2c_1, a_2c_1 - a_1c_2, a_1b_2 - b_1a_2).
    \end{equation}
    é chamado de \textbf{produto vetorial} de $\vec{u}$ por $\vec{v}$ e é denotado por $\vec{u}\times\vec{v}$.
\end{frame}

\begin{frame}
    \begin{proposicao}
        Sejam $\vec{u}$, \pause $\vec{v}$ \pause e $\vec{w}$ \pause vetores em $\real^3$ \pause e $\alpha \in \real$. \pause Temos\pause
        \begin{enumerate}[label=({\alph*})]
            \item $\vec{u}\times\vec{v} \pause = -\vec{v}\times\vec{u}$;\pause
            \item $\vec{u}\times\vec{v} = \vec{0}$ \pause se, e somente se, \pause $\vec{u}\varparallel\vec{v}$;\pause
            \item $\alpha(\vec{u}\times\vec{v}) \pause = (\alpha\vec{u})\times\vec{v} \pause = \vec{u}\times(\alpha\vec{v})$;\pause
            \item $\vec{u}\times(\vec{v} + \vec{w}) \pause = \vec{u}\times\vec{v} \pause + \vec{u}\times\vec{w}$.
        \end{enumerate}
    \end{proposicao}
\end{frame}

\begin{frame}
    \begin{observacao}
        O sentido de $\vec{u}\times\vec{v}$ \pause é dado pela regra da mão direita:\pause se o ângulo entre $\vec{u}$ e $\vec{v}$ \pause é $\theta$, \pause giramos o vetor $\vec{u}$ de um ângulo $\theta$ \pause até que coincida com o vetor $\vec{v}$ \pause e acompanhamos o movimento com os dedos da mão direita, \pause então o polegar vai apontar no sentido de $\vec{u}\times\vec{v}$.
    \end{observacao}
\end{frame}

\begin{frame}
    \begin{proposicao}\label{normaprodutovetorial}
        Sejam $\vec{u}$ e $\vec{v}$ vetores no espaço.\pause Então\pause
        \begin{equation}
            \norm{\vec{u}\times\vec{v}}^2 \pause = \norm{\vec{u}}^2\pause \norm{\vec{v}}^2\pause  - \inner{u}{v}^2.
        \end{equation}
    \end{proposicao}
\end{frame}

\begin{frame}
    \begin{proposicao}
        Quaisquer que sejam os vetores \pause não nulos \pause $\vec{u}$ \pause e $\vec{v}$ \pause em $\real^3$, \pause tem-se\pause
        \[
            \norm{\vec{u}\times\vec{v}}\pause  = \norm{\vec{u}}\norm{\vec{v}}\sin\theta
        \]
        onde $\theta$ \pause é o ângulo entre $\vec{u}$ e $\vec{v}$.
    \end{proposicao}
\end{frame}

\begin{frame}
    Agora considere o paralelogramo formado pelos vetores $\vec{u}$ e $\vec{v}$:
    \begin{figure}[!h]
        \centering
        %\caption{Interpretação geométrica do produto vetorial}
        \includegraphics{area-paralelogramo.pdf}
    \end{figure}
\end{frame}

\begin{frame}
    \begin{proposicao}
        Sejam $\vec{u}$, \pause $\vec{v}$ \pause e $\vec{w}$ \pause vetores em $\real^3$. \pause Então \pause
        \begin{enumerate}[label=({\alph*})]
            \item $\langle\vec{u}\times\vec{v}, \vec{w}\rangle \pause = \langle\vec{v}\times\vec{w}, \vec{u}\rangle$ \pause
            \item $\langle\vec{u}\times\vec{v}, \vec{w}\rangle \pause = \langle\vec{u}, \vec{v}\times\vec{w}\rangle$
        \end{enumerate}
    \end{proposicao}
\end{frame}

\begin{frame}
    Sejam $\vec{u}$, \pause $\vec{v}$ \pause e $\vec{w}$ \pause vetores em $\real^3$. \pause Considere o paralelepípedo definido por $\vec{u}$, \pause $\vec{v}$ \pause e $\vec{w}$:
    \begin{figure}[!h]
        \centering
        %\caption{Interpretação geométrica do produto misto}
        \includegraphics{volume-paralelepipedo.pdf}
    \end{figure}
\end{frame}

\end{document}
